\chapter{Mission and Subsystem Requirements}
\label{app:requirements}

This appendix reproduces the mission and subsystem requirements
specification. Requirement identifiers (MR- mission, FR- functional,
PR- performance, IF- interface, CR- constraint, SR- safety) are cited by
ID throughout the preceding chapters; the tables below give the full
requirement statement, verification method and acceptance criteria, and
owning subsystem for each. The full rationale for each requirement is
in the source specification, not reproduced here.

\begin{landscape}
\scriptsize

\begin{longtable}{p{1.4cm}p{8.2cm}p{10.2cm}p{1.9cm}}
\caption{Mission Requirements}
\label{tab:req-mission} \\
\toprule
\textbf{ID} & \textbf{Requirement} & \textbf{Verification \& Acceptance} & \textbf{Owner} \\
\midrule
\endfirsthead
\multicolumn{4}{l}{\tablename\ \thetable{} -- continued from previous page} \\
\toprule
\textbf{ID} & \textbf{Requirement} & \textbf{Verification \& Acceptance} & \textbf{Owner} \\
\midrule
\endhead
\midrule
\multicolumn{4}{r}{Continued on next page} \\
\endfoot
\bottomrule
\endlastfoot

MR-001 & The system shall autonomously detect and transition through six (6) flight phases: Standby, Launch, Ascent, Cruise, Descent, and Landing. & \textbf{Test.} Simulate each phase transition in laboratory environment; verify during balloon flight & FSW + CDH \\
MR-002 & The system shall operate continuously for a minimum of six (6) hours without external power input. & \textbf{Test.} Continuous operation endurance test with all subsystems active & EPS \\
MR-003 & The system shall acquire and log sensor data at a rate of at least one (1) Hz throughout the mission. & \textbf{Test.} Verify data logging rate during system operation; check for $\geq$99\% sample capture & CDH \\
MR-004 & The system shall transmit telemetry data to a ground station from altitudes up to forty (40) kilometers, corresponding to a worst-case slant range of 125~km assuming 120~km horizontal drift. & \textbf{Analysis + Test.} Link budget showing minimum 10~dB margin at 125~km slant range. & TT\&C \\
MR-005 & The system shall store all acquired flight data to non-volatile memory for post-flight retrieval. & \textbf{Test.} Verify 100\% of logged data is readable from SD card after power cycle & CDH \\
MR-006 & The system shall fit within a 1.5U CubeSat form factor of 100~mm $\times$ 100~mm $\times$ 170~mm. & \textbf{Examination.} Physical dimensional measurement of assembled system & Structure + All \\
MR-007 & The system shall survive environmental conditions including temperature range of $-60$\textdegree{}C to $+40$\textdegree{}C and pressure down to 1~kPa, with each subsystem operating within its own rated temperature range; payload-specific operating limits are captured in MR-010. & \textbf{Analysis + Test.} Component datasheet review for temperature/pressure ratings; thermal chamber test if available & All \\
MR-008 & The total flyable system mass shall not exceed 650 grams. & \textbf{Examination.} -- & All \\
MR-009 & The system shall host an on-board Edge-TPU ML payload that autonomously runs inference and delivers structured output to the OBC at 1~Hz within its own power and thermal envelope, without affecting the primary mission. & \textbf{Test + Analysis.} Both models (see payload PoC documentation) output to the OBC at 1~Hz; payload stays above 0\textdegree{}C with no throttle while other subsystems verify to their own ranges per MR-007; power within IF-009/PR-008; output verifiable against ground truth (FSM, SD frames) at documented accuracy. & Payload \\
\end{longtable}

\begin{longtable}{p{1.4cm}p{8.2cm}p{10.2cm}p{1.9cm}}
\caption{Functional Requirements}
\label{tab:req-functional} \\
\toprule
\textbf{ID} & \textbf{Requirement} & \textbf{Verification \& Acceptance} & \textbf{Owner} \\
\midrule
\endfirsthead
\multicolumn{4}{l}{\tablename\ \thetable{} -- continued from previous page} \\
\toprule
\textbf{ID} & \textbf{Requirement} & \textbf{Verification \& Acceptance} & \textbf{Owner} \\
\midrule
\endhead
\midrule
\multicolumn{4}{r}{Continued on next page} \\
\endfoot
\bottomrule
\endlastfoot

FR-001 & The OBC shall provide at least one (1) I2C buses for sensor interfaces. & \textbf{Examination.} Verify STM32 configuration shows $\geq$1 I2C peripherals enabled & OBC \\
FR-002 & The OBC shall provide at least one (1) SPI bus for high-speed peripheral communication. & \textbf{Examination.} Verify STM32 configuration shows $\geq$1 SPI peripheral enabled & OBC \\
FR-003 & The OBC shall provide at least one (1) UART interface for GPS communication. & \textbf{Examination.} Verify STM32 configuration shows $\geq$1 UART peripheral enabled & OBC \\
FR-004 & The FSW shall implement a finite state machine with six (6) distinct states corresponding to mission phases. & \textbf{Test.} Unit test verifying all six states are implemented and reachable & FSW \\
FR-005 & The FSW shall transition from Standby to Launch state when vertical acceleration exceeds 1.5g for at least two (2) consecutive seconds. & \textbf{Test.} Simulate acceleration input; verify transition occurs within 5 seconds of threshold & FSW \\
FR-006 & The FSW shall transition from Launch to Ascent state when altitude exceeds 100 meters AND vertical velocity exceeds 5~m/s. & \textbf{Test.} Simulate altitude and velocity inputs; verify both conditions required for transition & FSW \\
FR-007 & The FSW shall transition from Ascent to Cruise state when altitude rate of change is less than 0.5~m/s for at least thirty (30) consecutive seconds. & \textbf{Test.} Simulate altitude plateau; verify 30-second stability window required & FSW \\
FR-008 & The FSW shall transition from Cruise to Descent state when altitude rate of change is less than $-2$~m/s for at least ten (10) consecutive seconds. & \textbf{Test.} Simulate descent rate; verify 10-second confirmation window & FSW \\
FR-009 & The FSW shall transition from Descent to Landing state when GPS velocity is less than 1~m/s AND altitude is less than 200 meters for at least five (5) consecutive seconds. & \textbf{Test.} Simulate landing conditions; verify both conditions met for $\geq$5 seconds & FSW \\
FR-010 & The FSW shall persist the current flight phase state to non-volatile memory within one (1) second of any state transition. & \textbf{Test.} Power cycle system during each phase; verify correct phase restored after reboot in 6/6 test cases & FSW \\
FR-011 & The FSW shall implement a watchdog timer with a timeout period of ten (10) seconds. & \textbf{Test.} Fault injection test; verify system resets within 10 seconds when watchdog not serviced & FSW \\
FR-012 & The FSW shall prohibit transitions to undefined or invalid states. & \textbf{Test.} Attempt invalid state transitions; verify no undefined states reached during exhaustive testing & FSW \\
FR-013 & The CDH shall acquire GPS position data (latitude, longitude, altitude) at a rate of at least 1~Hz. & \textbf{Test.} Measure GPS data update rate; verify $\geq$1~Hz sustained output & CDH \\
FR-014 & The CDH shall acquire IMU data (3-axis acceleration, 3-axis rotation) at a rate of at least 1~Hz. & \textbf{Test.} Measure IMU data acquisition rate; verify $\geq$1~Hz output for all 6 axes & CDH \\
FR-015 & The CDH shall acquire barometric pressure and derived altitude data at a rate of at least 1~Hz. & \textbf{Test.} Measure barometer data acquisition rate; verify $\geq$1~Hz output & CDH \\
FR-016 & The CDH shall log to SD card: timestamp, current flight phase, GPS data, IMU data, barometer data, payload results, payload PoC verification footage and battery voltage. & \textbf{Examination.} Inspect SD card log file format; verify all specified fields present & CDH \\
FR-017 & The CDH shall write logged data to SD card at a rate of at least 1~Hz with $\geq$99\% write success rate. & \textbf{Test.} Verify $\geq$1 log entry per second during extended test; calculate write success percentage & CDH \\
FR-018 & The CDH shall implement equipment handlers interfacing with each sensor, handling startup, reading and flagging outdated data. & \textbf{Test.} Verify correct startup and reading in nominal mode and data flagging in fault mode & CDH \\
FR-019 & The CDH shall store the data obtained by the equipment handlers in the datapool vector accessible for FSW & \textbf{Test.} Verify correct interfacing with equipment handlers, availbillity and correct formatting for FSW & CDH \\
FR-020 & The CDH shall store a Spacecraft Configuration Vector (SCV) in a nonvolatile memory and update it regurlarly & \textbf{Test.} Verify correct storage and recovery capabillities & CDH \\
FR-021 & The CDH shall provide altitude data to FSW from GPS when available; otherwise from barometer. & \textbf{Test.} Block GPS signal; verify FSW receives barometer altitude and continues phase detection & CDH \\
FR-022 & The TT\&C shall transmit standarized telemetry packets at a rate of approximately one (1) packet per twenty (20) seconds, corresponding to a duty cycle below 1\% to comply with European ISM band regulations. & \textbf{Test.} Measure packet transmission rate over 5-minute interval & TT\&C \\
FR-023 & The telemetry packet shall include: header, packet type, sequence number, timestamp, current flight phase, GPS latitude, GPS longitude, GPS altitude, IMU acceleration (3-axis), selected payload data, and battery voltage, and parity check according to the ICD. & \textbf{Examination.} Capture and inspect transmitted packet structure; verify all fields present & TT\&C \\
FR-024 & The TT\&C shall operate at 868~MHz within the European ISM band (863--870~MHz) in compliance with ETSI EN 300 220. & \textbf{Examination.} Verify LoRa module hardware specification confirms 868~MHz operation & TT\&C \\
FR-025 & The ground station shall receive, decode and display standarised telemetry data in ICD in human-readable format. & \textbf{Demonstration.} Demonstrate telemetry reception and display during integration test & TT\&C \\
FR-026 & The TT\&C shall configure the LoRa radio module at Spreading Factor 9 (SF9), bandwidth 125~kHz, and coding rate 4/5, giving a data rate of 1758~bps and a receiver sensitivity of $-137$~dBm. & \textbf{Test.} Verify LoRa firmware register configuration shows SF9, BW 125~kHz, CR 4/5. Confirm data rate of 1758~bps and sensitivity of $-137$~dBm match RFM95W datasheet Table 3 and Table 5. & TT\&C \\
FR-027 & The OBC shall exchange data with the Coral Dev Board Micro payload over UART at 115200~baud, triggering inference and receiving a fixed 16-byte output block at 1~Hz. & \textbf{Test.} UART capture confirms 115200~baud and a correctly framed 16-byte block received at 1~Hz; the block is logged to SD and included in the outbound telemetry packet. & OBC \\
\end{longtable}

\begin{longtable}{p{1.4cm}p{8.2cm}p{10.2cm}p{1.9cm}}
\caption{Performance Requirements}
\label{tab:req-performance} \\
\toprule
\textbf{ID} & \textbf{Requirement} & \textbf{Verification \& Acceptance} & \textbf{Owner} \\
\midrule
\endfirsthead
\multicolumn{4}{l}{\tablename\ \thetable{} -- continued from previous page} \\
\toprule
\textbf{ID} & \textbf{Requirement} & \textbf{Verification \& Acceptance} & \textbf{Owner} \\
\midrule
\endhead
\midrule
\multicolumn{4}{r}{Continued on next page} \\
\endfoot
\bottomrule
\endlastfoot

PR-001 & Phase detection latency shall not exceed five (5) seconds from threshold condition being met to state transition completion. & \textbf{Test.} Measure time delay between threshold crossing and state transition in 95\% of test cases & FSW \\
PR-002 & Data logging shall capture at least 99\% of expected samples during normal operation. & \textbf{Test.} Calculate ratio: (actual samples logged) / (expected samples); verify $\geq$0.99 & CDH \\
PR-003 & GPS horizontal position accuracy shall be within ten (10) meters CEP (Circular Error Probable). & \textbf{Test.} Compare GPS output to known surveyed position during static ground test & CDH \\
PR-004 & The TT\&C system shall maintain a radio link with a minimum margin of 10~dB at a worst-case slant range of 125~km, corresponding to a balloon altitude of 40~km with 120~km horizontal drift. & \textbf{Analysis + Test.} Link budget calculation showing positive margin; ground range test at maximum feasible distance & TT\&C \\
PR-005 & The OBC shall operate at a clock frequency sufficient to run the FSW superloop within each 1~Hz cycle with margin. & \textbf{Examination.} Verify configured clock frequency in STM32 configuration $\geq$80~MHz & OBC \\
PR-006 & The system boot time from power-on to operational state shall not exceed thirty (30) seconds. & \textbf{Test.} Measure time from power application to first telemetry packet transmission & All \\
PR-007 & The TT\&C link budget shall demonstrate a minimum link margin of 10~dB and an SNR exceeding 0~dB at the worst-case slant range of 125~km. The effective SNR after LoRa processing gain shall exceed 10.5~dB to ensure BER below $10^{-6}$. & \textbf{Analysis.} Link budget showing link margin $>$10~dB, raw SNR $>$0~dB, and effective SNR after processing gain $>$10.5~dB at 125~km slant range. & TT\&C \\
PR-008 & The EPS power budget shall demonstrate a minimum margin of 1.5$\times$ between usable battery capacity and total energy required for the full flight duration of six (6) hours. & \textbf{Analysis.} Power budget calculation showing margin $\geq$1.5$\times$ at 6-hour flight duration with cold temperature derate factor of 0.9 applied to nominal battery capacity. & EPS \\
\end{longtable}

\begin{longtable}{p{1.4cm}p{8.2cm}p{10.2cm}p{1.9cm}}
\caption{Interface Requirements}
\label{tab:req-interface} \\
\toprule
\textbf{ID} & \textbf{Requirement} & \textbf{Verification \& Acceptance} & \textbf{Owner} \\
\midrule
\endfirsthead
\multicolumn{4}{l}{\tablename\ \thetable{} -- continued from previous page} \\
\toprule
\textbf{ID} & \textbf{Requirement} & \textbf{Verification \& Acceptance} & \textbf{Owner} \\
\midrule
\endhead
\midrule
\multicolumn{4}{r}{Continued on next page} \\
\endfoot
\bottomrule
\endlastfoot

IF-001 & The GPS module shall communicate with the OBC via UART at 9600~baud using NMEA 0183 protocol. & \textbf{Test.} Verify successful NMEA sentence parsing and position data extraction & CDH + OBC \\
IF-002 & The IMU shall communicate with the OBC via I2C at 100~kHz (standard mode) or 400~kHz (fast mode). & \textbf{Test.} Verify successful I2C communication and register read/write operations & CDH + OBC \\
IF-003 & The barometer shall communicate with the OBC via I2C at 100~kHz (standard mode) or 400~kHz (fast mode). & \textbf{Test.} Verify successful pressure and temperature data acquisition & CDH + OBC \\
IF-004 & The SD card shall communicate with the OBC via SPI at a clock frequency between 100~kHz and 25~MHz. & \textbf{Test.} Verify successful file creation, write, and read operations & CDH + OBC \\
IF-005 & The LoRa module shall communicate with the OBC via SPI at a clock frequency up to 10~MHz. & \textbf{Test.} Verify successful packet transmission and configuration register access & TT\&C + OBC \\
IF-006 & The CDH subsystem shall provide sensor data to FSW in a standardized struct format updated at 1~Hz. & \textbf{Examination.} Code review verifying data structure definition and update rate & CDH + FSW \\
IF-007 & The FSW shall provide current flight phase and system status to TT\&C in a defined packet format. & \textbf{Examination.} Code review and packet capture analysis & FSW + TT\&C \\
IF-008 & The EPS shall provide battery voltage monitoring output to the OBC via analog voltage divider and ADC. & \textbf{Test.} Verify voltage measurement accuracy within $\pm$5\% of actual battery voltage & EPS + OBC \\
IF-009 & The EPS shall supply a regulated 5V rail capable of sustaining a continuous current of at least 500~mA and peak transients of up to 2A to the Coral Dev Board Micro payload. & \textbf{Test.} Measure rail voltage during Coral inference burst; verify voltage remains above 4.5V throughout peak current transient. & EPS + OBC \\
\end{longtable}

\begin{longtable}{p{1.4cm}p{8.2cm}p{10.2cm}p{1.9cm}}
\caption{Constraint Requirements}
\label{tab:req-constraint} \\
\toprule
\textbf{ID} & \textbf{Requirement} & \textbf{Verification \& Acceptance} & \textbf{Owner} \\
\midrule
\endfirsthead
\multicolumn{4}{l}{\tablename\ \thetable{} -- continued from previous page} \\
\toprule
\textbf{ID} & \textbf{Requirement} & \textbf{Verification \& Acceptance} & \textbf{Owner} \\
\midrule
\endhead
\midrule
\multicolumn{4}{r}{Continued on next page} \\
\endfoot
\bottomrule
\endlastfoot

CR-001 & The system total mass shall not exceed 650 kilograms. & \textbf{Examination.} Weigh fully assembled flyable system on calibrated scale; verify total does not exceed 650~g. & Structure + All \\
CR-002 & The system shall use only COTS (Commercial Off-The-Shelf) components available within project budget. & \textbf{Examination.} Verify all components on approved BOM are commercially available & All \\
CR-003 & The system development shall be completed within twelve (12) weeks from project start. & \textbf{Examination.} Verify all subsystems integrated and tested by Week 11 per Gantt chart & All \\
CR-004 & The system total cost shall not exceed the approved project budget. & \textbf{Examination.} Compare actual expenditures to approved BOM budget & Project Mgmt \\
CR-005 & The system shall comply with German/EU regulations for radio transmission in the 868~MHz ISM band, including a maximum transmit power of 25~mW ERP ($+14$~dBm) and a duty cycle not exceeding 1\% in the 868.0--868.6~MHz sub-band per ETSI EN 300 220. & \textbf{Analysis.} Verify firmware sets transmit power within legal limit. Calculate actual TX duty cycle from packet size, data rate, and transmission interval; confirm below 1\%. & TT\&C \\
CR-006 & The GPS module shall support operation in Airborne $<$1g mode to avoid COCOM altitude limits. & \textbf{Examination + Test.} Verify GPS module datasheet specifies Airborne mode support; configure and test & CDH \\
\end{longtable}

\begin{longtable}{p{1.4cm}p{8.2cm}p{10.2cm}p{1.9cm}}
\caption{Safety Requirements}
\label{tab:req-safety} \\
\toprule
\textbf{ID} & \textbf{Requirement} & \textbf{Verification \& Acceptance} & \textbf{Owner} \\
\midrule
\endfirsthead
\multicolumn{4}{l}{\tablename\ \thetable{} -- continued from previous page} \\
\toprule
\textbf{ID} & \textbf{Requirement} & \textbf{Verification \& Acceptance} & \textbf{Owner} \\
\midrule
\endhead
\midrule
\multicolumn{4}{r}{Continued on next page} \\
\endfoot
\bottomrule
\endlastfoot

SR-001 & The system shall include over-current protection on all power distribution lines. & \textbf{Examination + Test.} Verify fuses or electronic current limiters present; test trip threshold & EPS \\
SR-002 & The system shall include reverse polarity protection on the main power input. & \textbf{Examination.} Verify diode or MOSFET reverse polarity protection circuit present & EPS \\
SR-003 & The system shall protect the battery against over-discharge by monitoring per-cell voltage via IF-008 and commanding an orderly firmware shutdown before any cell falls below its minimum safe voltage of approximately 3.0~V per cell, with hysteresis. & \textbf{Examination.} Hardware EN-pin UVLO backstop (2.8~V cut, 3.1~V recover per cell) verified independent of firmware state; IF-008 reading verified within plus or minus 5 percent of reference. The firmware-side orderly shutdown named in this requirement is not implemented -- per FMECA, no load shedding or power switching exists in firmware; accepted risk, mitigated pre-flight through battery sizing margin (see EPS chapter). & FSW + EPS \\
SR-004 & The system shall not pose a radio interference risk to balloon flight control systems. & \textbf{Test.} Coordinate with balloon operator; verify no interference during ground tests & TT\&C \\
SR-005 & If Iridium is added: the FSW shall not trigger Coral Dev Board Micro inference cycles simultaneously with Iridium transmission bursts. & \textbf{Test.} If Iridium is integrated, verify firmware schedules transmissions during Coral idle periods. Measure rail voltage during worst-case concurrent load scenario. & FSW + EPS \\
\end{longtable}

\end{landscape}
